\section{Related work}
\label{sec:related}

\paragraph{Agent benchmarks.}
CMCA adopts practices from established agent benchmarks.
From SWE-bench~\citep{jimenez2024swebench} it takes one submission folder per entry with mandatory trajectories.
From Terminal-Bench~\citep{merrill2026terminalbench} and its framework Harbor~\citep{harbor} it takes the task format (instruction, environment, oracle solution, verifier), containerized execution and pre-integrated CLI agents.
From $\tau$-bench~\citep{yao2024taubench} it takes the pass$^k$ reliability metric.
From the Holistic Agent Leaderboard~\citep{kapoor2025hal} it takes reporting cost next to accuracy and sharing all logs.
The Agentic Benchmark Checklist~\citep{zhu2025abc} supplies the task-validity, outcome-validity and reporting checks, including a null-agent baseline.
BIG-bench~\citep{srivastava2022bigbench} supplies the canary string in every task file.

\paragraph{Fault injection for agents.}
Several benchmarks inject tool failures into agent loops built for the study.
ReliabilityBench~\citep{gupta2026reliabilitybench} applies chaos-engineering fault profiles (timeouts, rate limits, partial responses, schema drift) to ReAct and Reflexion agents on two models.
ToolMaze~\citep{zhu2026toolmaze} crosses tool-graph complexity with explicit or implicit, transient or permanent perturbations, and finds implicit semantic failures the hardest.
AgentNoiseBench~\citep{wang2026agentnoisebench} injects user noise and tool noise into existing benchmarks.
ToolMisuseBench~\citep{sigdel2026toolmisusebench} uses deterministic simulators with declarative fault plans replayable under a seed.
PlanBench-XL~\citep{liu2026planbenchxl} blocks tools in a large retail tool ecosystem.
These works measure models and reasoning strategies in controlled environments; none of them runs a production coding-agent harness.

AgentChaosBench~\citep{zhang2026agentchaosbench} injects faults into five agentic systems to study a different question: diagnosing and localizing the fault from telemetry.
Xu and Chen~\citep{xu2026crossprovider} run fault-style probes (truncated input, cancellation, write attempts in read-only mode) against a cross-provider review contract with Claude and Codex backends; they test one contract, not a fault schedule over harness tool layers.

\paragraph{Harness effects.}
A growing body of work shows that the harness changes results when the model is fixed.
Harness-Bench~\citep{yao2026harnessbench} evaluates production harnesses, including Claude Code, Codex and OpenCode, across models on 106 tasks, and argues that capability should be reported per model and harness pair; it labels recovery after errors that occur naturally in trajectories but injects no faults.
Zhang et al.~\citep{zhang2026modelorharness} evaluate OpenCode, Codex CLI, Pi and DeepAgents with four models and report that harness rankings vary across models and benchmarks.
Hu et al.~\citep{hu2026beyondmodel} study harness components in software-engineering agents with open-weight models, and Lewis~\citep{lewis2026samemodel} compares two context-management configurations of one harness on coding benchmarks.
CMCA shares the view that the harness matters and adds a controlled fault schedule, so every configuration faces faults drawn from the same distribution, and on the service track the same bytes.

\paragraph{Developer tools.}
Open-source tools inject faults into agents their users build: at the model API (AgentChaos~\citep{agentchaos}), at the model API and MCP layers through a proxy (AgentBreak~\citep{agentbreak}, Saboteur~\citep{saboteur}), at recorded MCP tool responses (AgentCheck~\citep{mazumder2026agentcheck}), or around agent functions in Python (agentfuzz~\citep{agentfuzz}, chaos-agents~\citep{chaosagents}).
ChaosMonkey MCP~\citep{chaosmonkeymcp} reverses the roles: it gives a coding agent tools to fault-inject the code under development.

\paragraph{Position.}
To the best of our knowledge, as of November 2026 no prior work applies one seeded fault schedule inside released coding-agent harnesses, at both the MCP layer and the harnesses' built-in tool layer, and compares harness configurations under a fixed main model.
